\documentclass{reqengbook}

\title{
\bf\sffamily\fontsize{22}{28}\selectfont
Requirements Engineering for Software Developers
}

\author{\sffamily\fontsize{20}{30}\selectfont Björn Regnell}

\date{\vspace{2em}\sffamily\small Updated: \today 
\\ License: CC-BY-SA 
\\ \url{https://github.com/bjornregnell/reqeng-book} 
}

\begin{document}
%\pagenumbering{roman} 

\begin{minipage}{0.5\textwidth}
  \maketitle%
\end{minipage}%
\begin{minipage}{0.5\textwidth}
  \centering\vspace{2.2cm}
  \includegraphics[width=0.9\textwidth]{../img/phone-support.jpg}%
\end{minipage}%
\pagebreak

\setcounter{tocdepth}{2}

\begin{multicols*}{3}  % * means unbalanced columns allowed
\fontsize{8.5}{11}\selectfont

\tableofcontents
\end{multicols*}

\newgeometry{hmargin={45mm,148.5mm},vmargin={15mm,15mm}}

\part{Intentions}

%!TEX encoding = UTF-8 Unicode
%!TEX root = ../book/reqeng-book.tex

\chapter{Context}%

\section{Context of RE}

\subsection{Business model}
Funding model. Economy of dev.

License fee. Freemium. Subscription.

\subsection{Customer-supplier relationship}

\subsection{Product and integration}

Type of product

\begin{itemize}
  \item Level of customization (customer specific, generic, mixed)
  \item Hardware integration (Pure SW, embedded HW+SW, pure HW)
  \item Network integration (stand-alone, network connection, cloud: back-end -- front-end)
\end{itemize}

\subsection{Service and delivery}

\url{https://en.wikipedia.org/wiki/Product_lifecycle}

\url{https://en.wikipedia.org/wiki/Software_as_a_service}

Delivery model: local executable, cloud service, on-prem self-hosting, saas, etc etc


\section{Domain modeling}

\subsection{Stakeholder modeling}

\subsection{Goal modeling}

Helps or Hurts.

Related to both elicitation (stakeholder intentions) and specification (user roles and goals).

\TODO{Decide how much on Goals here or more details under Elicitation}???

\subsection{Context diagram}

Inner and outer domain.

\TODO{Decide if it should be called 'Context Diagram' instead?}






%!TEX encoding = UTF-8 Unicode
%!TEX root = ../book/reqeng-book.tex

\chapter{Elicitation}%
blabla


Elicitation Goals \& Challenges%

\section{What information to elicit?}

\begin{itemize}
  \item Present
  \item Future 
  \item Consensus 
  \item Commitment 
  \item Innovative ideas
  \item Goals -- Domain -- Product -- Design
  \item Priorities (cost, benefit, risk of requriements)
\end{itemize}

\section{Why is elicitation so hard?}

Elicitation challenges (barriers/difficulties) 

\section{Interacting with Stakeholders}%

\subsection{Surveys}

\subsection{Interviews}

blabla

\subsection{Case Studies}

Demonstrations

Observation of Current work 

Prototyping -> its own chapter 

Usability testing -> own subchapter in Module on Quality

Pilot system deployment

\subsection{Requirements from Operation}

support issues, telemetry, user forums and social media, A/B-testing experiments


\section{Business Intelligence}
blabla  (sv Omvärldsanalys)

study similar companies, competitors, subcontractors, suppliers 

\section{Creativity and Innovation}

trigger ideas 

assess novelty and market 

brainstorming

focus-groups

%!TEX encoding = UTF-8 Unicode
%!TEX root = ../book/reqeng-book.tex

\chapter{Prioritization}%
blabla

\section{Why prioritize?}%

\section{The prioritization process}

\subsection{Activities}
who does what when 

\subsection{Challenges}

Finding a good abstraction level

Combinatorial explosion

Inter-dependencies

Not easy to predict the future

Power and politics

\subsection{Aspects}

Example aspects:

-- Importance (e.g. financial benefit, urgency, strategic value, market share...)

--  Penalty (e.g. bad-will if requirement not included)

--  Cost (e.g., staff effort)

--  Time (e.g., lead time)

--  Risk (e.g., technical risk, business risk)

--  Volatility (e.g. scope instability, probability of change)

-- Other aspects: competitors, brand fitness, competence, release theme...

-- Combination of aspects: cost vs. benefit, cost vs. risk, importance vs. volatility

-- Optimize: minimize or maximize some combinations, e.g., cost vs. benefit

\section{Scales}

\section{Methods}

\textbullet Grouping, Numerical assignment (grading)

-- Can be done using any scale (categorical, ordinal, ratio)

-- Quick and easy; but a risk is that all reqs are deemed highly important as they are
not challenged against each other; may be misinterpreted as ratio scale (even if
”4” not necessarily is "twice as much" as ”2” when using an ordinal scale).

\textbullet 100-dollar-test

-- Ratio scale, quick and easy, risk of shrewd tactics (listigt taktikspel)

\textbullet Ranking (sorting)

-- Ordinal scale, pairwise comparison, easy and rather quick

\textbullet Top-ten (or Top-n)

-- Ordinal scale if the top list is ranked / Categorical grouping if not ranked; very
quick and simple, gives a rough estimate on a limited set of req

\textbullet (Analytical Hierarchy Process (AHP))

-- Ratio scale, pairwise comparison, tool is needed, redundancy gives measure of
consistency 


\part{Functionality}

\chapter{Requirements Engineering Tools}

\section{Natural Language}

-- controlled domain-terminology 

-- explicit or implicit reqs, 

-- shall / should / must / optional MOSCOW EARS

-- point to RM chapter on textual feature requirements

\section{Generic tools}
  word processors, presentation tools, drawing tools , issue trackers 

  AI and Large language models 

\section{RE-specific tools}
  Features of RE Tools: Store, Collaborate, Document Model, Language, traceability, different views, RE-aware AI support, ...

  Examples: Doors, RequisitePro, 

\subsection{reqT}

more details in appendix

%!TEX encoding = UTF-8 Unicode
%!TEX root = ../book/reqeng-book.tex

\chapter{Data Requirements}%
blabla


\section{Data Dictionary}
blabla

\section{Entity-Relationship Diagrams}
blabla

\section{Class Diagrams}
blabla

\section{Data Screens}
blabla  (aka Virtual Windows)

\section{Formal Data Specification}%

\subsection{Regular Expressions}

\subsection{Protocol Buffers}

\url{https://en.wikipedia.org/wiki/Protocol_Buffers}

%!TEX encoding = UTF-8 Unicode
%!TEX root = ../book/reqeng-book.tex

\chapter{Logic}%

Functional Requirements on business logic that provide constraints on system behavior.


\section{When to decide user-system work split?}

\section{Contextual usage modeling}%

\subsection{User stories}

A user story is an informal, general explanation of a software feature written from the perspective of the end user. Its purpose is to articulate how a software feature will provide value to the customer.

\url{https://en.wikipedia.org/wiki/User_story}

\url{https://www.atlassian.com/agile/project-management/user-stories}

\subsection{Use cases and tasks}

\subsection{Narratives and personas}

Narrative: Rich and detailed description of a real-world usage situation in the form of a short-story.

Persona: a fictional character representing an imagined human user, often attributed with concrete personal attitudes and desires. 


\section{System behavior modeling}

-- Domain Events

-- Product Events

-- Product Functions triggered by Events: (Input, State) => Output

\subsection{State diagram}%
\url{https://en.wikipedia.org/wiki/State_diagram} 

\subsection{Interaction diagrams}

\begin{itemize}
  \item UML Sequence Diagrams \url{https://en.wikipedia.org/wiki/Sequence_diagram}
  \item ITU Message Sequence Charts \url{https://en.wikipedia.org/wiki/Message_sequence_chart}
\end{itemize}

\subsection{Data-flow diagrams}%
\url{https://en.wikipedia.org/wiki/Data-flow_diagram} 





\chapter{Prototyping}

\section{Why prototype?}

\section{Media}

\subsection{Screens and mockups}


\chapter{Delegated requirements}

\section{Standards as requirements}

\subsection{Regulatory requirements}

GDPR 

Point to chapter on AI-act in "Special Domains"

\section{Prototypes as Requirements}

\section{Test Cases as Requirements}

\section{Development Process Requirements}


\part{Quality}

\chapter{Quality Requirements}

\chapter{Requirements Validation}

\chapter{Product Scoping}

\section{Epics and Themes}

bundles of features and use cases that deliver value

\section{Release Planning} 

\chapter{Product Verification}
\section{System Testing}
\section{Acceptance Testing}
\section{Regression Testing}

\part{Special Contexts}

\chapter{Agile RE}

\chapter{Open Source RE}

\chapter{Continuous RE}

\section{CI/CD}
\section{DevOps}

\chapter{RE for Artificial Intelligens}

\chapter{RE for High-Assurance Systems}

\chapter{RE in Hardware-Software Co-Design}

\part{Appendicies}\appendix

\chapter{Requirement Modelling with reqT}

\chapter{Lab 1: Prioritization}

\chapter{Lab 2: Release Planning}

\end{document}
